\silentchapter{Anexe}

\section{Comenzi de rulare utilizate}
\label{anexa:comenzi}

Această anexă prezintă comenzile principale folosite pentru instalarea și testarea aplicației.

\begin{code}
\begin{minted}{bash}
pip install -r requirements.txt

python main.py --mode backtest --pair EURUSD --start 2018-01-01 \
  --end 2024-12-31 --timeframe D

python main.py --mode backtest --pair GBPUSD --start 2018-01-01 \
  --end 2024-12-31 --timeframe D

python main.py --mode backtest --pair USDJPY --start 2018-01-01 \
  --end 2024-12-31 --timeframe D
\end{minted}
\caption{Comenzi pentru instalare și backtesting}
\label{code:anexa_backtest_commands}
\end{code}

Pentru modul live paper trading, aplicația se pornește după ce TWS sau IB Gateway rulează în cont paper și are activat accesul API:

\begin{code}
\begin{minted}{bash}
python main.py --mode live --pairs EURUSD,GBPUSD,USDJPY --timeframe 15m \
  --ib-port 7497 --client-id 14
\end{minted}
\caption{Comandă pentru rularea în mod live paper}
\label{code:anexa_live_command}
\end{code}

Pentru monitorizarea live se pornesc serviciile Prometheus și Grafana prin Docker Compose:

\begin{code}
\begin{minted}{bash}
docker compose up -d
\end{minted}
\caption{Pornirea serviciilor de observabilitate}
\label{code:anexa_observability_start}
\end{code}

Adresele utilizate pentru verificare sunt:
\begin{itemize}
    \item \texttt{http://localhost:18000/metrics} pentru metricile expuse de bot;
    \item \texttt{http://localhost:39090/targets} pentru starea țintei Prometheus;
    \item \texttt{http://localhost:33000} pentru interfața Grafana.
\end{itemize}

\section{Metrici de observabilitate live}
\label{anexa:metrici_prometheus}

Pentru monitorizarea modului live paper, aplicația expune metrici Prometheus care descriu starea curentă a botului, prețurile urmărite și pozițiile asociate perechilor valutare. Aceste metrici sunt utile atât pentru verificarea rapidă a funcționării sistemului, cât și pentru construirea dashboard-urilor Grafana.

\begin{table}[H]
    \centering
    \small
    \begin{tabular}{|p{5.2cm}|p{8.7cm}|}
        \hline
        \textbf{Metrică} & \textbf{Rol} \\
        \hline
        \texttt{bot\_live\_price\{pair\}} & Prețul live pentru fiecare pereche valutară monitorizată. \\
        \hline
        \texttt{bot\_pair\_signal\{pair\}} & Semnalul numeric generat de allocator pentru perechea respectivă. \\
        \hline
        \begin{tabular}[t]{@{}l@{}}\texttt{bot\_pair\_}\\\texttt{open\_position\{pair\}}\end{tabular} & Indicator binar care arată dacă perechea are o poziție deschisă. \\
        \hline
        \begin{tabular}[t]{@{}l@{}}\texttt{bot\_pair\_}\\\texttt{position\_size\_units\{pair\}}\end{tabular} & Dimensiunea poziției, exprimată în unități ale instrumentului tranzacționat. \\
        \hline
        \begin{tabular}[t]{@{}l@{}}\texttt{bot\_pair\_}\\\texttt{realized\_pnl\_usd\{pair\}}\end{tabular} & Profitul sau pierderea realizată pentru perechea valutară. \\
        \hline
        \begin{tabular}[t]{@{}l@{}}\texttt{bot\_pair\_}\\\texttt{unrealized\_pnl\_usd\{pair\}}\end{tabular} & Profitul sau pierderea nerealizată pentru poziția curentă. \\
        \hline
        \texttt{bot\_equity\_usd} & Capitalul curent raportat de aplicație. \\
        \hline
        \texttt{bot\_open\_positions} & Numărul total al pozițiilor deschise. \\
        \hline
    \end{tabular}
    \caption{Metrici Prometheus expuse de aplicația live}
    \label{tabel:anexa_metrici_prometheus}
\end{table}

Exemple de interogări Prometheus utile pentru verificare sunt:

\begin{code}
\begin{minted}{bash}
up{job="forex_quant_bot"}
bot_live_price{pair="EURUSD"}
bot_pair_signal{pair="USDJPY"}
bot_pair_position_size_units{pair="USDJPY"}
bot_pair_realized_pnl_usd{pair="USDJPY"}
\end{minted}
\caption{Interogări Prometheus pentru monitorizarea botului}
\label{code:anexa_prometheus_queries}
\end{code}

\section{Procedură de verificare pentru rularea live paper}
\label{anexa:procedura_live_paper}

Înaintea pornirii modului live paper, este necesară verificarea mediului de execuție și a conexiunilor externe. Procedura recomandată este:

\begin{enumerate}
    \item pornirea TWS sau IB Gateway în cont paper;
    \item verificarea portului API, de regulă \texttt{7497} pentru TWS paper sau \texttt{4002} pentru IB Gateway paper;
    \item pornirea botului cu perechile valutare dorite și cu un \texttt{client-id} unic;
    \item verificarea endpoint-ului \texttt{/metrics} în browser;
    \item verificarea stării \texttt{UP} în pagina \texttt{/targets} din Prometheus;
    \item deschiderea dashboard-ului Grafana corespunzător perechii valutare monitorizate;
    \item verificarea în Interactive Brokers a ordinelor paper generate și a pozițiilor curente.
\end{enumerate}

Această procedură reduce riscul de configurare greșită și permite confirmarea faptului că botul, brokerul paper și sistemul de observabilitate funcționează împreună.

\section{Structura artefactelor generate}
\label{anexa:artefacte}

Fiecare rulare creează un director separat în \texttt{output}. Numele directorului include modul de rulare, perechea valutară, timeframe-ul, perioada testată și timestamp-ul rulării.

\begin{table}[H]
    \centering
    \small
    \begin{tabular}{|p{4.2cm}|p{9.4cm}|}
        \hline
        \textbf{Fișier} & \textbf{Rol} \\
        \hline
        \texttt{trades.csv} & Tranzacțiile închise, cu prețuri de intrare/ieșire, P\&L, regim și motivul semnalului. \\
        \hline
        \texttt{transactions.csv} & Evenimentele individuale de intrare și ieșire pentru fiecare tranzacție. \\
        \hline
        \texttt{equity\_curve.csv} & Evoluția capitalului după fiecare bară procesată. \\
        \hline
        \texttt{returns.csv} & Tabel sumar pentru profit brut, pierdere brută, profit factor și comisioane. \\
        \hline
        \texttt{strategy\_metrics.csv} & Metrici asociate performanței strategiilor. \\
        \hline
        \texttt{regime\_metrics.csv} & Performanța grupată după regimul de piață al intrării. \\
        \hline
        \begin{tabular}[t]{@{}l@{}}\texttt{risk\_overlay\_}\\\texttt{metrics.csv}\end{tabular} & Deciziile componentei de risc și motivele aprobării sau respingerii semnalelor. \\
        \hline
        \texttt{summary.txt} & Raport text cu indicatorii principali și lista tranzacțiilor. \\
        \hline
        \texttt{dashboard.html} & Dashboard HTML pentru inspectarea vizuală a performanței. \\
        \hline
    \end{tabular}
    \caption{Fișiere generate de o rulare de backtesting}
    \label{tabel:anexa_artefacte}
\end{table}

\section{Parametri relevanți de configurare}
\label{anexa:config}

Aplicația folosește clase de configurare pentru strategii, risc, date, logare și broker. Parametrii principali ai componentei de risc sunt:

\begin{table}[H]
    \centering
    \small
    \begin{tabular}{|p{4.8cm}|p{3cm}|p{6cm}|}
        \hline
        \textbf{Parametru} & \textbf{Valoare implicită} & \textbf{Descriere} \\
        \hline
        \texttt{initial\_capital} & 10.000 USD & Capitalul inițial al simulării. \\
        \hline
        \texttt{risk\_per\_trade} & 0,01 & Ponderea capitalului riscată per tranzacție. \\
        \hline
        \texttt{max\_total\_exposure} & 0,05 & Expunerea totală maximă permisă. \\
        \hline
        \texttt{drawdown\_circuit\_breaker} & 0,05 & Pragul de drawdown care oprește temporar tranzacționarea. \\
        \hline
        \texttt{atr\_stop\_multiplier} & 1,5 & Multiplicatorul ATR pentru stop-loss. \\
        \hline
        \begin{tabular}[t]{@{}l@{}}\texttt{take\_profit\_}\\\texttt{atr\_multiplier}\end{tabular} & 5,0 & Multiplicatorul ATR pentru take-profit. \\
        \hline
        \texttt{max\_notional\_leverage} & 1,0 & Levierul noțional maxim permis. \\
        \hline
    \end{tabular}
    \caption{Parametri de risc utilizați în proiect}
    \label{tabel:anexa_risc}
\end{table}

\clearpage
\section{Fragment reprezentativ de flux decizional}
\label{anexa:flux}

Fragmentul următor sintetizează logica folosită de motorul de backtesting. Codul complet se află în directorul sursă al aplicației.

\begin{code}
\begin{minted}[basicstyle=\footnotesize\ttfamily\bfseries\color{codeIdentifier}]{python}
for event in BarReplay(market_data, warmup_bars=warmup_bars):
    bar = event.bar
    close_price = float(bar["close"])

    decisions = {
        strategy.name: strategy.evaluate(event.history)
        for strategy in strategies
    }
    composite = allocator.allocate(
        decisions,
        current_regime=str(bar["regime"]),
        performance_tracker=tracker,
    )

    risk_decision = risk_overlay.evaluate_entry(
        pair=pair,
        open_positions=[],
        bias=composite.bias,
        equity=equity,
        current_open_risk_ratio=open_risk_ratio,
        close_price=close_price,
        atr_value=float(bar["atr"]),
        quote_to_usd_rate=quote_to_usd,
    )
\end{minted}
\caption{Sinteză a fluxului de decizie din backtesting}
\label{code:anexa_flux_decizie}
\end{code}

\section{Fragment reprezentativ pentru dimensionarea poziției}
\label{anexa:position_size}

Dimensionarea poziției se realizează pornind de la riscul acceptat și de la distanța stop-loss-ului exprimată prin ATR.

\begin{code}
\begin{minted}[basicstyle=\footnotesize\ttfamily\bfseries\color{codeIdentifier}]{python}
target_risk_amount = equity * self.config.risk_per_trade
stop_distance = atr_value * self.config.atr_stop_multiplier
usd_risk_per_unit = stop_distance * quote_to_usd_rate

size_by_risk = int(max(target_risk_amount / usd_risk_per_unit, 1))
max_notional_usd = equity * self.config.max_notional_leverage
size_by_notional = int(max(max_notional_usd / usd_notional_per_unit, 0))

size_units = min(size_by_risk, size_by_notional)
\end{minted}
\caption{Dimensionarea poziției în funcție de risc și volatilitate}
\label{code:anexa_position_size}
\end{code}
